\chapter{Counting processes, spike trains, and information}
A spike train is a set of event times or, more generally, a random
locally finite counting measure. Write $N((a,b])$ for the number of
events in an interval. A conditional intensity $\lambda(t\mid
\mathcal H_t)$ describes the conditional instantaneous event rate
given the preceding history $\mathcal H_t$, if such a predictable
intensity exists.

\entry{71}{Kolmogorov extension theorem}
{A consistent collection of finite-dimensional probability distributions
defines a stochastic process on a suitable product measurable space.}
{Build process models from their joint laws across finite time sets.}
\entry{72}{Campbell's theorem}
{For a point process with mean measure $\Lambda$ and nonnegative
measurable $h$, $\E[\sum_{t_i} h(t_i)]=\int h\,d\Lambda$;
signed integrable variants require absolute integrability.}
{Compute expected spike-triggered sums and event statistics.}
\entry{73}{Poisson superposition and thinning}
{Independent Poisson processes superpose to a Poisson process with
the sum intensity; independent thinning with valid retention
probability yields the corresponding Poisson intensity.}
{Generate spike-train simulation benchmarks and control models.}
\entry{74}{Doob--Meyer decomposition}
{A suitable integrable counting submartingale admits a
predictable compensator $A(t)$, and $N(t)-A(t)$ is a martingale;
when absolutely continuous, $A(t)=\int_0^t\lambda(s)\,ds$.}
{Connect conditional intensity estimation and martingale residuals.}
\entry{75}{Time-rescaling theorem}
{Under a correctly specified continuous conditional intensity and
regularity of event times, transformed inter-event compensator
increments are iid $\operatorname{Exp}(1)$ (with standard caveats
for boundaries and conditioning).}
{Construct goodness-of-fit diagnostics for spike-train GLMs.}
\entry{76}{Conditional point-process likelihood}
{For a simple point process with conditional intensity $\lambda$
on a finite observation interval, under standard domination
assumptions its log-likelihood, up to reference terms, is
$\sum_i\log\lambda(t_i)-\int\lambda(t)\,dt$.}
{Derive maximum-likelihood fitting for spiking models.}
\entry{77}{Palm theory and Slivnyak's theorem}
{For homogeneous Poisson processes, a reduced Palm process at a
specified event has the law of the original process; more
general Palm laws need not have this property.}
{Understand event-centered descriptions and null spike correlations.}
\entry{78}{Gibbs' inequality}
{For probability measures $P,Q$ with $P\ll Q$,
$\KL(P\|Q)\geq0$, with equality if $P=Q$ almost everywhere.}
{Understand KL-based divergence estimators and likelihood objectives.}
\entry{79}{Data-processing inequality}
{For Markov chain $X\to Y\to Z$,
$I(X;Z)\leq I(X;Y)$ whenever mutual information is defined.}
{Characterize information lost through feature extraction or decoding.}
\entry{80}{Fano's inequality}
{For a discrete class variable with $K\geq2$ classes,
entropy and classification error obey
$H(X\mid\widehat X)\leq h_2(P_e)+P_e\log(K-1)$,
with a consistent logarithm base.}
{Relate decoder error to information; do not confuse with the
spike-count Fano \emph{factor}.}

\begin{example}[Count likelihood and binning]
With constant intensity $\lambda$ on interval length $\Delta$,
$N_\Delta\sim\Poisson(\lambda\Delta)$.
Then
$\log p(n\mid\lambda)=n\log(\lambda\Delta)
-\lambda\Delta-\log(n!)$.
A binned GLM approximates continuous-time spiking only relative
to its chosen bin width and assumptions about multiple spikes per bin.
\end{example}
\begin{caution}
A time-rescaling test diagnoses discrepancies from the hypothesized
conditional intensity; a passed test is not proof of a unique
physiological mechanism. Model misspecification can also arise from
latent common drive or refractory dynamics.
\end{caution}
